\subsubsection{Provádění skoku a nastavení fyzikálních parametrů}

Metoda \texttt{Jump()} se volá, když jsou splněny všechny podmínky pro skok (dostatečný coyote time a jump buffer). Nejprve se nastaví \texttt{isJumping} na true, což indikuje, že hráč je aktuálně ve vzduchu a nelze skočit znovu. Současně se nastaví \texttt{jumpCut} na false, což resetuje stav přerušení skoku.

Následně se vynuluje vertikální rychlost na \texttt{Rigidbody2D} a poté se nastaví na hodnotu \texttt{jumpForce} z \texttt{PlayerControllerData}. Toto zajistí konzistentní počáteční rychlost skoku bez ohledu na předchozí pohyb. Důležité je také vynulování časových proměnných \texttt{lastOnGroundTime} a \texttt{lastPressedJumpTime}, aby se skok neopakoval v dalších snímcích.

Hodnota \texttt{jumpForce} se typicky pohybuje v rozmezí 10-20 v závislosti na nastavení gravitace a požadované výšce skoku. Správné nastavení této hodnoty je klíčové pro celkový herní pocit.

\begin{lstlisting}[language=csharp,caption={Listing 7: Metoda Jump pro provedení skoku}]
private void Jump()
{
    if (rb == null || data == null) return;

    isJumping = true;
    jumpCut = false;
    lastPressedJumpTime = 0;
    lastOnGroundTime = 0;

    rb.linearVelocity = new Vector2(rb.linearVelocity.x, 0);
    rb.linearVelocity = new Vector2(rb.linearVelocity.x, data.jumpForce);
}
\end{lstlisting}
